\section{Decode an image in stages}
A PNG file does not contain a row of framebuffer words. It contains chunks, some of
which carry parts of a compressed stream. Decompression produces rows of filtered
sample bytes. Reversing the row filters recovers red, green and blue samples. Only
then can the program pack those samples into the words used by the bitmap display.
This chapter builds these transformations separately. When a pixel is wrong, inspect
the bytes at the preceding boundary before changing the display loop.

The first image is two pixels wide and five rows high. Each row uses a different
filter. Its first IDAT chunk contains only the first byte of the zlib header; the next
chunk contains the rest of the stream. This deliberate split prevents a convenient
file layout from becoming an accidental decoder assumption. The fixture helper at
\path{drafts/png-fixtures/fixture.py} reproduces the file and each intermediate byte
sequence. It generates inputs and expected observations; it does not implement a
PNG decoder for the assembly program to call.

The reading begins with scalar processing. Later, sixteen independent byte additions
can replace part of the Up-filter loop, and vector shuffles can pack RGB samples.
The scalar version remains a comparison throughout. A faster image with even one
incorrect channel is a failed decoder.
